
\documentclass{article}
\usepackage{colortbl}
\usepackage{makecell}
\usepackage{multirow}
\usepackage{supertabular}

\begin{document}

\newcounter{utterance}

\centering \large Interaction Transcript for game `wordle\_withclue', experiment `high\_frequency\_words\_clue\_no\_critic', episode 4 with reflexion{-}Qwen3.5{-}4B.
\vspace{24pt}

{ \footnotesize  \setcounter{utterance}{1}
\setlength{\tabcolsep}{0pt}
\begin{supertabular}{c@{$\;$}|p{.15\linewidth}@{}p{.15\linewidth}p{.15\linewidth}p{.15\linewidth}p{.15\linewidth}p{.15\linewidth}}
    \# & \multicolumn{2}{c}{Player} && \multicolumn{2}{c}{Game Master} \\
    \hline

    \theutterance \stepcounter{utterance}  
    & & & \multicolumn{4}{p{0.6\linewidth}}{
        \cellcolor[rgb]{0.9,0.9,0.9}{
            \makecell[{{p{\linewidth}}}]{
                \texttt{\tiny{[P1$\langle$GM]}}
                \texttt{You have attempted this same task before and did not succeed. These are the plans you wrote after those attempts.} \\
\texttt{Trial \#1: New plan: The last attempt was aborted because the explanation text became excessively long, rambling, and repetitive, likely violating the constraint to provide a "concise explanation." I failed to synthesize the clue and the feedback into a single, direct sentence before generating the guess. Next time, I will strictly limit the explanation to one or two sentences that directly state the confirmed letters (e.g., 'e' is at index 1, 'r' is present but not at index 3) and the resulting pattern, then immediately provide the guess without any backtracking or speculation about the clue's meaning. I will ensure the total response length is minimal and adheres strictly to the "explanation:" and "guess:" format.} \\
\\ 
\texttt{Trial \#2: New plan: The last attempt was aborted because the explanation became excessively long, repetitive, and speculative, violating the constraint for a "concise explanation." I failed to synthesize the clue and the feedback into a single, direct sentence before generating the guess. Next time, I will strictly limit the explanation to one or two sentences that directly state the confirmed letters (e.g., 'e' is at index 1, 'r' is present but not at index 3) and the resulting pattern, then immediately provide the guess without any backtracking or speculation about the clue's meaning. I will ensure the total response length is minimal and adheres strictly to the "explanation:" and "guess:" format.} \\
\\ 
\\ 
\texttt{You are a language wizard who likes to guess words by using the given rules.} \\
\\ 
\texttt{Welcome to Wordle! You have six attempts to guess the target word, a valid English word of five lowercase letters (a{-}z). Please use the tags "explanation:" and "guess:" to provide a concise explanation for each guess.} \\
\\ 
\texttt{To help you make an informed guess, you will receive a clue for the word, such as} \\
\texttt{clue: snowy white.} \\
\\ 
\texttt{Here is an example guess based on the clue:} \\
\texttt{explanation: In the fairy tail Snow White, the girl is killed because she eats a poisoned apple. And the word apple has 5 letters.} \\
\texttt{guess: apple} \\
\\ 
\texttt{After each guess, your answer will be validated, and you will receive feedback indicating which letters are correct (green), which letters are correct but in the wrong position (yellow), and which letters are incorrect (red). This feedback can be useful in determining which letters to include or exclude in your next guess.} \\
\\ 
\texttt{For example, the feedback for "apple" might be:} \\
\texttt{guess\_feedback: a<yellow> p<yellow> p<green> l<yellow> e<red>} \\
\\ 
\texttt{The explanation should contain details about how both the guess\_feedback and clue work together in tandem to formulate a new guess.} \\
\\ 
\texttt{Do not generate any other text.} \\
\\ 
\texttt{Let's begin with your first guess.} \\
\\ 
\texttt{clue: 1985 hit by heart} \\
            }
        }
    }
     \\ \\

    \theutterance \stepcounter{utterance}  
    & \multicolumn{4}{p{0.6\linewidth}}{
        \cellcolor[rgb]{0.561,0.737,0.561}{
            \makecell[{{p{\linewidth}}}]{
                \texttt{\tiny{[P1$\rangle$GM]}}
                \texttt{explanation: The clue "1985 hit by heart" refers to the song "My Heart Will Go On" by Celine Dion, but that is too long. A better fit is "When You Say Nothing At All" or perhaps a specific song title. Wait, the most famous 1985 hit related to the heart is likely "The Power of Love" by Huey Lewis and the News, or maybe "I Want to Know What Love Is" by Foreigner? Actually, a very famous 1985 hit with "heart" in the title or theme is "Heart" by Boney M? No. Let's reconsider. "1985 hit by heart" might refer to the song "Heart" by Boney M? No. How about "Heart of Glass"? No, that's 1980. Maybe "I Love Rock 'n' Roll"? No. Let's try a common 5{-}letter word. "Love" is 4. "Song" is 4. "Hit" is 3. Maybe the word is "Rock"? No. Let's try "Pop". No. Wait, could it be "Beat"? "Beat" is a 1985 hit? "Beat It" is 1982. "Heart" is 1985? Yes, "Heart" by Boney M is 1985. But "Heart" is 5 letters. Let's guess "heart".} \\
\texttt{guess: heart} \\
            }
        }
    }
    & & \\ \\

    \theutterance \stepcounter{utterance}  
    & & & \multicolumn{2}{p{0.3\linewidth}}{
        \cellcolor[rgb]{0.9,0.9,0.9}{
            \makecell[{{p{\linewidth}}}]{
                \texttt{\tiny{[GM$|$GM]}}
                \texttt{attempts = 1} \\
\texttt{target\_word = never} \\
\texttt{target\_word\_clue = 1985 hit by heart} \\
\texttt{guess = heart} \\
\texttt{guess\_feedback = h<red> e<green> a<red> r<yellow> t<red>} \\
            }
        }
    }
    & & \\ \\

    \theutterance \stepcounter{utterance}  
    & & & \multicolumn{4}{p{0.6\linewidth}}{
        \cellcolor[rgb]{0.9,0.9,0.9}{
            \makecell[{{p{\linewidth}}}]{
                \texttt{\tiny{[P1$\langle$GM]}}
                \texttt{guess\_feedback: h<red> e<green> a<red> r<yellow> t<red>} \\
\\ 
\texttt{Provide your response only in this format.} \\
\texttt{explanation: details} \\
\texttt{guess: word} \\
            }
        }
    }
     \\ \\

    \theutterance \stepcounter{utterance}  
    & \multicolumn{4}{p{0.6\linewidth}}{
        \cellcolor[rgb]{0.561,0.737,0.561}{
            \makecell[{{p{\linewidth}}}]{
                \texttt{\tiny{[P1$\rangle$GM]}}
                \texttt{explanation: The feedback shows 'e' is correct at index 1, 'r' is present but at index 3, while 'h', 'a', and 't' are incorrect. The word pattern is \_e\_r\_.} \\
\texttt{guess: ever} \\
            }
        }
    }
    & & \\ \\

    \theutterance \stepcounter{utterance}  
    & & & \multicolumn{2}{p{0.3\linewidth}}{
        \cellcolor[rgb]{0.9,0.9,0.9}{
            \makecell[{{p{\linewidth}}}]{
                \texttt{\tiny{[GM$|$GM]}}
                \texttt{Error: The length of the guessed word is not 5.} \\
            }
        }
    }
    & & \\ \\

    \theutterance \stepcounter{utterance}  
    & & & \multicolumn{2}{p{0.3\linewidth}}{
        \cellcolor[rgb]{0.9,0.9,0.9}{
            \makecell[{{p{\linewidth}}}]{
                \texttt{\tiny{[GM$|$GM]}}
                \texttt{game\_result = ABORT} \\
            }
        }
    }
    & & \\ \\

    \theutterance \stepcounter{utterance}  
    & & \multicolumn{4}{p{0.6\linewidth}}{
        \cellcolor[rgb]{0.95,0.95,0.95}{
            \makecell[{{p{\linewidth}}}]{
                \texttt{\tiny{[P1$\leftrightarrow$P]}}
                \texttt{Post{-}episode Reflexion (not sent to GM):} \\
\texttt{New plan: The last attempt was aborted because the explanation was too verbose and speculative, failing to meet the "concise" requirement. I spent too many sentences guessing song titles and analyzing the clue instead of immediately processing the feedback. Next time, I will strictly limit the explanation to one sentence that directly states the confirmed letters (e.g., 'e' is at index 1, 'r' is present but not at index 3), deduces the resulting pattern (\_e\_r\_), and proposes a new guess without any backtracking or speculation about the clue's meaning. I will ensure the response is minimal and adheres strictly to the "explanation:" and "guess:" format.} \\
            }
        }
    }
    & \\ \\

\end{supertabular}
}

\end{document}
